%

% %%%

%%%   Самарский государственный технический университет

%%%   Кафедра прикладной математики и информатики

%%%

%%%   Преамбула для статьи в журнал "Вестник Самарского государственного

%%%   технического университета. Серия: «Физико-математические науки»"

%%%   Ориентирована под MikTEx 2.4 for Win32

%%%

%%%   Хотя сам журнал собирается под GNU/Linux на дистрибутиве TeXLive



\documentclass[11pt,a4paper,twoside]{article}



\usepackage[cp1251]{inputenc}

\usepackage[T2A]{fontenc}

\usepackage[english,russian]{babel}

\usepackage{samgtubib}

\usepackage[dvips]{graphicx}

\usepackage[multiple]{footmisc}



\usepackage{samgtu}

\renewcommand{\VestnikYear}{2009} % Год издания Вестника

\renewcommand{\VestnikNumber}{2\,(19)} % Номер Вестника

\renewcommand{\VestnikMonth}{Ноябрь} % Месяц выхода Вестника

\renewcommand{\VestnikData}{01.11.09} % Дата подписания Вестника в печать

\renewcommand{\VestnikUPL}{26,64} % Усл. печ. л.

\renewcommand{\VestnikUIL}{25,73} % Уч.-изд. л.

\renewcommand{\VestnikRegNumber}{479/09} % Регистрационный номер

\renewcommand{\VestnikOrder}{994} % Номер заказа







%%

% \begin{document}



\renewcommand{\firstpage}{361}

\renewcommand{\lastpage}{368}



%%%%%%%%%%%%%%%%%%%%%%%%%%%%%%%%%%%%%%%%%%%%%%%%%%%%%%%%%%%%%%%%%%%%%%%%%%%%%%%%%%%%

%Информация о статье и об авторах на русском и английском языках



% Номер УДК

\udk{514.84}%Геометрические методы в квантовой механике и в теории элементарных частиц

\msc{51P05, 81Q70}% Коды MSC см. http://www.ams.org/msc/

%51P05 - Geometry and physics

%81Q70 - Differential-geometric methods, including holonomy, Berry and Hannay phases, etc.



\title{Кривизна фазового пространства}% Полный заголовок

{Кривизна фазового пространства}% Заголовок, который идёт в верхний колонтитул



\author{М.\,Г.~Иванов}% Автор(ы) для заголовка, если авторов несколько укать \inst

{М.~Г.~И\,в\,а\,н\,о\,в}% Автор(ы) для оглавления и колонтитула через \,



\institute{Московский физико-технический институт (государственный университет),\\

Россия, 141700, Московская область, Долгопрудный, Институтский переулок, 9.}





\email{E-mail}{ivanov.mg@mipt.ru} %E-mail's авторов



\otherinf{\fullauthor{Михаил Геннадьевич Иванов}\regalia{к.ф.-м.н., доц.} \position{доцент}

\dept{каф. теоретической физики} }



\keywords{фазовое пространство, квантовая механика, калибровочная симметрия, кривизна, некоммутативная геометрия}



\workabstract{

Показывается, что электромагнитное поле в классической и квантовой

механике естественным образом описывается через геометрию

расширенного фазового пространства, в число координат которого

входят время и сопряжённый ему импульс $p_0=-E$.}



\engtitle{Phase space curvature}



\engauthor{M.\,G.~Ivanov}{M.~G.~I\,v\,a\,n\,o\,v}





\enginstitute{Moscow Institute of Physics and Technology (State University),\\

9, Inststitutskii per., Dolgoprudnyi, Moskovskaya obl.,  141700, Russia.}% Место работы каждого из авторов на англ. языке

% Если несколько, то так  \enginstitute{ Организация 1 \and Организация 2  \and Организация 3}



%Moscow Institute of Physics and Technology  Dolgoprudnyi, Moskovskaya obl., Russia

%Address: 	,   Inststitutskii per., 9 



\otherenginf{\fullauthor{Mikhail G. Ivanov} \regalia{Ph.\,D. (Phys. \& Math.)} \position{Associate Professor}

\dept{Dept. of Theoretical Physics}}



\engabstract{Electromagnetic field in classical and quantum mechanics is naturally represented

by geometry of extended phase space, with extra coordinates of time and canonically conjugate momentum

$p_0=-E$.}



\engkeywords{phase space, quantum mechanics,

gauge symmetry, curvature, noncommutative geometry}



\date{17/XII/2012} % Дата поступления статьи в редакцию.

\revisiondate{15/II/2013} % В окончательном виде



%%%%%%%%%%%%%%%%%%%%%%%%%%%%%%%%%%%%%%%%%%%%%%%%%%%%%%%%%%%%%%%%%%%%%%%%%%%%%%%%%%%%





\maketitle





\Section[n]{Введение}

 Квантовое коммутационное соотношение между координатой и~сопряжённым импульсом имеет вид $[\hat Q, \hat P]=i\hbar$,  оно часто переписывается через операторы конечных сдвигов по координате и импульсу \cite{ivanovmg:bib:Weyl,ivanovmg:bib:Faddeev}:

\begin{equation}\label{Uu}

  \hat U(u)=e^{-\frac{iu\hat P}\hbar},\qquad

  \hat U(u)\psi(Q)=\psi(Q-u),\qquad

  \hat U(u)\psi(P)=e^{-\frac{iuP}\hbar}\psi(P),

\end{equation}

\begin{equation}\label{Vv}

  \hat V(v)=e^{-\frac{iv\hat Q}\hbar},\qquad

  \hat V(v)\psi(Q)=e^{-\frac{ivQ}\hbar}\psi(Q),\qquad

  \hat V(v)\psi(P)=\psi(P+v).

\end{equation}

  Коммутационное соотношение записывается в

  виде группового коммутатора

\begin{equation}\label{UVUV}

  \hat U(u)\hat V(v)\hat U^{-1}(u)\hat V^{-1}(v)=

  \hat U(u)\hat V(v)\hat U(-u)\hat V(-v)=

  e^{\frac{iuv}\hbar}.

\end{equation}



  Таким образом, последовательность сдвигов

  в фазовой плоскости $(Q,P)$ по замкнутому прямоугольному

  контуру площади $uv$ соответствует умножению волновой функции на

  фазовый множитель $e^{\frac{iuv}\hbar}$. Это означает, что

  фазовой плоскости можно приписать постоянную кривизну в

  расслоении над группой $U(1)$ аналогично кривизне, задаваемой

  в калибровочной $U(1)$ теории тензором электромагнитного поля.



  Можно связать между собой две

  хорошо разработанные области математической физики:

  симплектическую геометрию и теорию калибровочных полей.

  Симплектическая форма и тензор электромагнитного поля

  объединяются в один объект, задающий кривизну в расслоении

  фазового пространства над группой $U(1)$.

  В литературе эта аналогия разрабатывается в одну сторону:

  квантовая механика как калибровочная теория в фазовом

  пространстве \cite{ivanovmg:bib:deGosson}.



Для частицы во внешнем электромагнитном поле характеристики поля уже не должны входить в

 гамильтониан, вместо этого в духе общей теории относительности (ОТО) заряженная частица

 свободно движется в искривлённом фазовом пространстве.



  Хотя исходная идея связана с квантовыми коммутационными соотношениями,

  подобная переформулировка возможна как для классической, так и для квантовой механики.





\Section{Коммутаторы и скобки Пуассона}

  Пусть $X^K$ --- координаты в фазовом пространстве.

  Для коммутаторов и классических скобок Пуассона имеем

$$

  [\hat X^K,\hat X^L]=i\hbar J^{KL},\qquad

  \{X^K,X^L\}=J^{KL}.

$$

  В нашей интерпретации $J^{KL}$ --- тензор кривизны

 фазового пространства (над группой $U(1)$).



  В канонических координатах

$$

  X^{Q^i}=Q^i,\qquad X^{P_j}=P_j,\qquad

  J^{Q^iP_j}=-J^{P_jQ^i}=\delta^i_j,\quad

  J^{Q^iQ^j}=J^{P_iP_j}=0.

$$

 Переход от обобщённых импульсов $P$ к кинематическим импульсам $p$

 позволяет исключить статическое магнитное поле из гамильтониана,

 описав его как добавку к кривизне фазового пространства.

   Для таких (<<новых канонических>>) координат $x$ имеем

\begin{equation}\label{x-ncanon}

  x^{q^i}=q^i=Q^i,\qquad

  x^{p_j}=p_j=P_j-\frac{e}cA_j(Q),\qquad

  \{x^K,x^L\}=I^{KL},

\end{equation}

\begin{equation}\label{I-ncanon}

  I^{q^ip_j}=-I^{p_jq^i}=\delta^i_j,\qquad

  I^{q^iq^j}=0,\qquad

  I^{p_ip_j}=\frac{e}c F_{ij}=\frac{e}c(\partial_i A_j-\partial_j A_i).

\end{equation}

 Симплектическая форма $\omega$ задаётся матрицей, обратной

 к матрице $I$, т.~е. $\omega_{KL}I^{LM}=\delta_K^M$.

\begin{equation}\label{s-form-ncanon}

  \omega_{q^ip_j}=-\omega_{p_jq^i}=-\delta^i_j,\qquad

  \omega_{q^iq^j}=\frac{e}c F_{ij},\qquad

  \omega_{p_ip_j}=0.

\end{equation}



  В случае статических (не зависящих от времени) полей переход к

  <<новым каноническим>> координатам --- это всего лишь замена

  координат в фазовом пространстве, $I^{KL}$ --- это прежний

  тензор $J^{KL}$ в новых координатах, гамильтониан --- прежняя

  скалярная функции на фазовом пространстве, скобка Пуассона также

  остаётся прежней. Это автоматически доказывает (в силу произвольности

  выбора координат при тензорной записи скобки Пуассона),

  что и динамика системы остаётся прежней.



   Поскольку между старыми и новыми координатами есть

   взаимно однозначное соответствие, мы будем нумеровать их

   одинаковыми индексами.

%\vspace{3mm}



\smallskip



\begin{newthm}{Утверждение 1} Движение классической или квантовой заряженной частицы во внешнем статическом магнитном поле $\mathbf{H}$ может быть описано гамильтонианом свободной частицы $$H=\frac{\mathbf{p}^2}{2m},$$ если на фазовом пространстве

задана симплектическая форма вида \eqref{s-form-ncanon}{\rm ,} включающая в себя магнитное поле $F_{ij}=-e_{ijk}H_k.$

Скобки Пуассона $($классические или квантовые$)$ задаются неканоническими соотношениями \eqref{I-ncanon}{\rm .}

Это описание соответствует гамильтониану $$H=\frac{(\mathbf{P}-\frac{e}{c}\mathbf{A})^2}{2m}$$ с

векторным потенциалом $\mathbf{A}$ $($где $\mathbf{H}=\mathrm{rot}\,\mathbf{A}),$ не зависящим от времени{\rm ,} в неканонических

координатах $x$ \eqref{x-ncanon}{\rm .}

\end{newthm}



\smallskip



Похожий подход в рамках некоммутативной геометрии развивался Беллиссардом \cite{ivanovmg:bib:bellissard}.

%\vspace{3mm}



\smallskip



  Чтобы в рамках данного подхода описать  переменное

  электромагнитное поле, необходимо расширить фазовое пространство,

  рассматривая время $x^0=t$ и соответствующий времени обобщённый импульс

  $p_0=-E$ как дополнительные координаты.



\Section{Лагранжев формализм в расширенном конфигурационном пространстве}

Пусть задан некий лагранжиан $L(t,q,\dot q)$, который явно зависит

от времени. Мы можем переписать неавтономную систему как автономную,

записав время $t$ как функцию от некоторого монотонного параметра

$\tau$:

$$

  S[q(t)]=\int L(t,q,\dot q)\,dt

  =\int L(t,q,q'/t')\,t'\,d\tau,\qquad

\dot q=\frac{dq}{dt},\quad q'=\frac{dq}{d\tau}.

$$

Мы получаем новый (\emph{расширенный}) функционал действия

$S_\textrm{р}[t(\tau),q(\tau)]$ с новым (\emph{расширенным})

лагранжианом $L_\textrm{р}$, который для любой траектории совпадает

с исходным действием $S[q(t)]$, но зависит от другого набора функций:

$$

  S_\textrm{р}[t(\tau),q(\tau)]=\int

  L_\textrm{р}(t,t',q,q')\,d\tau,\qquad

  L_\textrm{р}(t,t',q,q')=L(t,q,q'/t')\,t'.

$$

 Легко проверить, что уравнения Эйлера для старого и расширенного

 лагранжианов эквивалентны, причём

$$

  p_i=\frac{\partial L_\textrm{р}}{\partial q^{i\prime}}

  =\frac{\partial L}{\partial\dot q^i},\qquad

  p_0=\frac{\partial L_\textrm{р}}{\partial t'}

  =L-\frac{\partial L}{\partial\dot

  q^i}\,\dot q^i=-E.

$$



Уравнения Лагранжа для расширенного лагранжиана для координат

$q(\tau)$ и времени $t(\tau)$ с точностью до множителя $t'$

совпадают с уравнениями Лагранжа и уравнением баланса энергии для

исходного лагранжиана

$$

  \frac{\delta S_\textrm{р}}{\delta q^i(\tau)}

 =t'\,\frac{\delta S}{\delta q^i(t)}

 =t'\Bigl(\frac{\partial L}{\partial q^i}-\frac{d p_i}{dt}\Bigr),\qquad

  \frac{\delta S_\textrm{р}}{\delta t(\tau)}

 =t'\Bigl(\frac{\partial L}{\partial t}+\frac{dE}{dt}\Bigr).

$$



<<Энергия>> для расширенного лагранжиана тождественно равна нулю:

$$

  \mathcal{E}=p_0t'+p_iq^{i\prime}-L_\textrm{р}=t'(p_0+E)\equiv0.

$$



Расширенное действие описывает ту же физическую

систему, что и исходное, но поскольку выбор параметризации времени

$t(\tau)$ произволен (любая гладкая монотонная функция),  

уравнения Лагранжа оказываются зависимыми (уравнение

баланса энергии выражается через остальные уравнения).



\smallskip



\Section{Гамильтонов формализм в расширенном фазовом пространстве}

Для расширенного фазового пространства перейдём к гамильтоновому

формализму. Соответствующее преобразование Лежандра оказывается

неоднозначным. <<Энергию>> для расширенного лагранжиана надо

выразить через координаты и импульсы (включая $t$ и $p_0$):

$$

 \mathcal{E}=t'p_0+q^{i\prime}p_i-L_\textrm{р}

 =t'(p_0+p_i\dot q^i-L)

 =t'\,(p_0+H(t,p,q)),

$$

  $t'$ не может быть определено из уравнений

  Лагранжа. Положим $$t'=f(t,p_0,q,p),$$

 где $f\not=0$ --- произвольная гладкая функция.

 Получается <<гамильтониан>>

\begin{equation}\label{H=f(P+H)}

  \mathcal{H}(t,p_0,q,p)=f(t,p_0,q,p)\cdot(p_0+H(t,q,p)).

\end{equation}

  Соответствующие уравнения Гамильтона (<<расширенные уравнения Гамильтона>>) имеют вид

%\begin{eqnarray*}

%\frac{dt}{d\tau}&=&\frac{\partial\mathcal{H}}{\partial p_0}

%=f+\frac{\partial f}{\partial p_0}(p_0+H),\\

%\frac{dp_0}{d\tau}&=&-\frac{\partial\mathcal{H}}{\partial t}

%=-f\,\frac{\partial H}{\partial t}-\frac{\partial f}{\partial t}(p_0+H),\\

%\frac{dq^i}{d\tau}&=&\frac{\partial\mathcal{H}}{\partial p_i}

%=f\,\frac{\partial H}{\partial p_i}+\frac{\partial f}{\partial p_i}(p_0+H),\\

%\frac{dp_i}{d\tau}&=&-\frac{\partial\mathcal{H}}{\partial q^i} =-f\,\frac{\partial

%H}{\partial q^i}-\frac{\partial f}{\partial q^i}(p_0+H).

%\end{eqnarray*}

\begin{eqnarray*}

\frac{dt}{d\tau}=\frac{\partial\mathcal{H}}{\partial p_0}

=f+\frac{\partial f}{\partial p_0}(p_0+H),~~

\frac{dp_0}{d\tau}=-\frac{\partial\mathcal{H}}{\partial t}

=-f\,\frac{\partial H}{\partial t}-\frac{\partial f}{\partial t}(p_0+H),\\

\frac{dq^i}{d\tau}=\frac{\partial\mathcal{H}}{\partial p_i}

=f\,\frac{\partial H}{\partial p_i}+\frac{\partial f}{\partial p_i}(p_0+H),~~

\frac{dp_i}{d\tau}=-\frac{\partial\mathcal{H}}{\partial q^i} =-f\,\frac{\partial

H}{\partial q^i}-\frac{\partial f}{\partial q^i}(p_0+H).

\end{eqnarray*}



<<На энергетической поверхности>>, т.~е. если задать начальные условия, для которых

 $p_0=-H$, расширенные уравнения Гамильтона дают уравнение хода

 времени $\frac{dt}{d\tau}=f$, уравнение баланса энергии и  уравнения Гамильтона для исходных координат и импульсов с новым  временем $\tau$. При этом воспроизводится гамильтонова динамика

 исходной системы.



 Если $f=\textrm{const}\not=0$, то для любых начальных условий

 воспроизводится гамильтонова динамика исходной системы, при этом

$$\frac{d~}{dt}(p_0+H)=0,$$ т.~е. начальное значение $p_0$ может

 быть произвольным и энергия $E=-p_0$ оказывается определена с

 точностью до произвольного постоянного слагаемого. Можно сказать, что

 в этом случае интеграл движения $p_0+H(t,q,p)$ --- нулевой уровень

 энергии, который может быть выставлен произвольно.

%\vspace{3mm}



\smallskip



\begin{newthm}{Утверждение 2} Движение классической заряженной частицы во внешнем электромагнитном поле $F_{ij}$ может быть описано 

<<гамильтонианом>> \eqref{H=f(P+H)} свободной частицы $$\mathcal{H}=f\cdot(p_0+\frac{\mathbf{p}^2}{2m}),$$ если на \emph{расширенном} фазовом пространстве (включающем время и энергию) задана симплектическая форма  \eqref{s-form-ncanon}{\rm ,} включающая электромагнитное поле $F_{ij}.$

Описание соответствует <<гамильтониану>> 

$$\mathcal{H}=f\cdot\Bigl(P_0-\frac{e}{c}A_0+\frac{(\mathbf{P}-\frac{e}{c}\mathbf{A})^2}{2m}\Bigr)$$ с

потенциалом $A_i=(A_0,\mathbf{A})$

% (где $F_{ij}=\partial_i A_j-\partial_j A_i$),

 в неканонических

координатах $x$ \eqref{x-ncanon} на расширенном фазовом пространстве{\rm .}

\end{newthm}

%\vspace{3mm}



\smallskip



\Section{Квантовая механика с расширенным гамильтонианом}

  В квантовом случае импульс в <<координатном по времени>> представлении

  $\hat p_0=-i\hbar\frac{\partial~}{\partial t}$.

  При $f=1$ имеем 

\begin{equation}\label{qHq}

\mathcal{\hat H}=\hat p_0+\hat H,

\end{equation} 

что даёт обычное уравнение Шрёдингера:

\begin{equation}\label{Hpsi=0}

  \mathcal{\hat H}\psi=0\quad\Leftrightarrow\quad

  \Bigl(-i\hbar\frac{\partial~}{\partial t}+\hat H\Bigr)\psi(t)=0.

\end{equation}



  При ином выборе функции $f\not=\mathrm{const}$ мы можем также

  воспроизвести уравнение Клейна"--~Фока"--~Гордона:

$$

  \mathcal{H}=f\cdot(p_0+H),\qquad f=p_0-H,\qquad

  H=\sqrt{m^2c^4+\mathbf{p}^2c^2},

$$

$$

  \mathcal{H}=p_0^2-H^2=p_0^2-\mathbf{p}^2 c^2-m^2c^4,\qquad

  \mathcal{\hat H}%=\hat p_0^2-\mathbf{\hat p}^2 c^2-m^2c^4

  =c^2\hbar^2\Bigl(\triangle-\frac1{c^2}\frac{\partial~}{\partial t^2}\Bigr)-m^2c^4.

$$



  Вернёмся, однако, к случаю $f=1$ как простейшему и, вероятно, наиболее фундаментальному.



  В энергетическом (импульсном по времени) представлении

  $p_0=-E$, и мы получаем

  стандартное стационарное уравнение Шрёдингера

$$

  \bigl(-E+\hat H\bigr)\psi(E)=0

$$

 с нормировкой на вероятность данной энергии

$\langle\psi(E)|\psi(E)\rangle=p_E$

 и нестандартной интерпретацией: $\psi(t)$ и $\psi(E)$ связаны

 между собой преобразованием Фурье.



  Несмотря на то, что в расширенный классический гамильтониан

время входит на общих основаниях с другими координатами, в квантовой механике между ними возникает различие за счёт определения

пространства волновых функций, которые не являются квадратично

интегрируемыми по времени.



  Уравнение \eqref{Hpsi=0} имеет вид уравнения на собственную функцию оператора $\mathcal{\hat H}$ с собственным числом 0.

  Мы можем написать аналогичное уравнение для произвольного собственного числа $E_0\in\mathbb{R}$:

\begin{equation}\label{Hpsi=E0}

  \mathcal{\hat H}\psi_{E_0}=E_0\quad\Leftrightarrow\quad

  \Bigl(-i\hbar\frac{\partial~}{\partial t}+\hat H\Bigr)\psi_{E_0}(t)=E_0.

\end{equation}

  Уравнения \eqref{Hpsi=E0} описывает ту же самую динамику, что и уравнение \eqref{Hpsi=0}, но с нулевым уровнем

  энергии, сдвинутым на $E_0$. Решения этих уравнений отличаются на фазовый множитель

\begin{equation}\label{psi0psiE0}

  \psi_{E_0}(t)=e^{i\frac{E_0}{\hbar}t}\psi(t).

\end{equation}

 Ненормируемость волновых функций $\psi_{E_0}$ при интегрировании по всем координатам,

 включая время, связана с тем, что функции относятся к непрерывному спектру.

 В стандартной формулировке квантовой механики, где время рассматривается не как координата,

 а как параметр, естественное пространство волновых функций --- гильбертово пространство $L_2(\mathbb{R}^n)$,

 при включении времени в число координат волновая функция должна принадлежать оснащённому расширенному гильбертовому

 пространству $\Phi'\supset L_2(\mathbb{R}^{n+1})\supset\Phi$.

%\vspace{3mm}



\smallskip



\begin{newthm}{Утверждение 3} Движение квантовой заряженной частицы во внешнем электромагнитном поле $F_{ij}$ может быть описано 

волновой функцией из оснащённого расширенного гильбертова пространства $\Phi'\supset L_2(\mathbb{R}^4)\supset\Phi,$

которая является собственной функцией непрерывного спектра <<гамильтониана>> \eqref{qHq} свободной частицы $$\mathcal{\hat H}=\hat p_0+\frac{\mathbf{\hat p}^2}{2m}$$ для произвольного собственного числа{\rm .}

Квантовые скобки Пуассона задаются неканоническими соотношениями \eqref{I-ncanon}{\rm ,} включающими

электромагнитное поле $F_{ij}.$

Это описание соответствует стандартному временному уравнению Шрёдингера с гамильтонианом

%\begin{equation}\label{Ht-old}

$$

\hat H=-\frac{e}{c}A_0+\frac{(\mathbf{\hat P}-\frac{e}{c}\mathbf{A})^2}{2m},

$$

%\end{equation} 

с потенциалом $A_i=(A_0,\mathbf{A})$ $($где $F_{ij}=\partial_i A_j-\partial_j A_i)$ в неканонических

координатах на расширенном фазовом пространстве $x$ \eqref{x-ncanon} и с волновой функцией

$\psi(t,\mathbf{r})$ которая рассматривается не как функция $t\to L_2(\mathbb{R}^3),$ а как элемент

 пространства 

$\Phi'\supset L_2(\mathbb{R}^4)\supset\Phi.$

\end{newthm}

%\vspace{3mm}



\smallskip



\Section{Пространство волновых функций}

 Однако пока волновые функции "--- по-прежнему функции на конфигурационном пространстве, хотя и расширенном.

 Чтобы получить функции, зависящие как от координат (включая~$t$), так и от импульсов (включая $p_0=-E$), воспользуемся

 представлением волновых функций как элементов оснащённого пространства Фока.



 Далее будем считать, что координаты и импульсы обезразмерены.



 Когерентное состояние $\psi_\mathbf{z}$ (где $\mathbf{z}=\frac{\mathbf{q}_0+i\mathbf{p}_0}{\sqrt2}\in\mathbb{C}^n$, $\mathbf{q}_0,\mathbf{p}_0\in\mathbb{R}^n$)

 получается из основного состояния $n$-мерного изотропного гармонического осциллятора

$$

  \psi_0(\mathbf{q})=\frac1{\pi^{n/4}}e^{-\frac{\mathbf{q}^2}2}

$$

 с помощью сдвига по координате на $\mathbf{q}_0$ и импульсу на $\mathbf{p}_0$, т.~е. с помощью некоторой

 комбинации операторов \eqref{Uu} и \eqref{Vv} для разных координат и импульсов.

 В зависимости от того, вдоль какой траектории в фазовом пространстве осуществляется сдвиг,

 когерентные состояния с одинаковым $\mathbf z$ могут отличаться друг от друга на фазовый множитель,

 что связано с некоммутативностью сдвигов по координате и импульсу \eqref{UVUV}.



 Набор когерентных состояний $\{\psi_\mathbf{z}\}_{\mathbf{z}\in\mathbb{C}^n}$ не является базисом

 (поскольку переопределён), но образует ортоподобную систему:

\begin{equation}\label{ortho}

\int_{\mathbb{C}^n}|\psi_\mathbf{z}\rangle\langle\psi_\mathbf{z}|\,d\mathbf{z}\,d\mathbf{z}^*

=\pi^n\cdot\hat 1.

\end{equation}

 Любой волновой функции $\psi(\mathbf{q})$ мы можем сопоставить функцию $\Psi(\mathbf{q}_0,\mathbf{p}_0)=\Psi(\mathbf{z})$:

$$

  \Psi(\mathbf{z})=\langle\psi_{\mathbf{z}^*}|\psi\rangle.

$$

 В силу \eqref{ortho} функция $\Psi$ обладает многими свойствами настоящей волновой функции, например,

$$

  \langle\psi_1|\psi_2\rangle=\frac1{\pi^n}\int \Psi_1^*\Psi_2\,d\mathbf{z}\,d\mathbf{z}^*.

$$

 $|\Psi|^2$ в классическом пределе переходит в совместное распределение вероятности на фазовом пространстве.



 Поскольку когерентные состояния определены с точностью до фазового множителя, функция $\Psi(\mathbf{q},\mathbf{p})$ определена

 с точностью до умножения на $e^{i\alpha(\mathbf{q},\mathbf{p})}$, т.~е. с точностью до локального калибровочного

 предобразования группы $U(1)$.



 Если все когерентные состояния были получены из $\psi_0$ сдвигом вдоль прямой в фазовом пространстве

 (это условие фиксирует калибровку), то

\begin{equation}\label{fok-gauge}

  \Psi(\mathbf{z})=e^{-\frac{\mathbf{zz}^*}2}\,f(\mathbf{z}),

\end{equation}

 где $f(\mathbf{z})$ --- функция из пространства Фока.



  Пространство Фока \cite{ivanovmg:bib:Faddeev,ivanovmg:bib:ivanovmg} (в указанных книгах пространство Фока

 описывается, но термин <<пространство Фока>> не используется) --- пространство аналитических функций $f(\mathbf{z})$ комплексных переменных $\mathbf{z}=(z_1,\dots,z_n)$, на котором определено скалярное произведение

\begin{equation}\label{f1f2}

 (f_1,f_2)=\frac1{\pi^n}\int e^{-\mathbf{zz}^*}\, f_1^*f_2\,d\mathbf{z}\,d\mathbf{z}^*.

\end{equation}

 В пространство Фока $\mathcal{D}$ включаются только те функции, для которых скалярный квадрат $(f,f)$ конечен.



 На пространстве Фока

\begin{equation}\label{aaqp}

\hat a_k=\frac{\partial~}{\partial z_k},\quad \hat a_k^+=z_k,\quad

  \hat q_k=\frac{\frac{\partial~}{\partial z_k}+z_k}{\sqrt2},\quad

  \hat p_k=\frac{\frac{\partial~}{\partial z_k}-z_k}{i\sqrt2},

\end{equation}

 что позволяет построить представление других операторов, выражающихся через координаты и импульсы.



 Мы можем определить расширенное пространство Фока, добавив ещё одну комплексную переменную $z_0$.



 Пространство Фока является гильбертовым пространством, и мы можем определить оснащённое пространство

 Фока $F'\supset \mathcal{D}\supset F$, выделив в $\mathcal{D}$ линейное всюду плотное подпространство $F$.



 Таким образом, в качестве аналога волновых функций на фазовом пространстве мы можем взять \eqref{fok-gauge}:

\begin{equation}\label{Psi(qp)}

 \Psi({q},{p})=e^{-\frac{q^2+p^2}4}\,f\Bigl(\frac{{q}-i{p}}{\sqrt2}\Bigr),\quad

 q=(q_0,q_1,\dots,q_n),\quad

 p=(p_0,p_1,\dots,p_n),

\end{equation}

где функция $f:\mathbb{C}^{n+1}\to\mathbb{C}$ принадлежит оснащённому расширенному пространству Фока.

Соответствующее представление гамильтониана может быть построено с помощью представления \eqref{aaqp}

операторов координаты и импульса, полагая $z=\frac{q-ip}{\sqrt2}$.

%\vspace{3mm}



\smallskip



\begin{newthm}{Утверждение 4} Собственная функция из оснащённого расширенного гильбертова пространства

$\Phi'\supset L_2(\mathbb{R}^4)\supset\Phi$ $($см{\rm .} Утверждение $3)$

для <<гамильтониана>>

$$

\mathcal{\hat H}=\hat P_0-\frac{e}{c}\hat A_0+\frac{(\mathbf{\hat P}-\frac{e}{c}\mathbf{\hat A})^2}{2m}

$$

 может быть представлена как элемент $f$ оснащённого расширенного пространства Фока $F'\supset \mathcal{D}\supset F,$ на котором <<гамильтониан>> \eqref{qHq}

действует согласно правилу \eqref{aaqp}{\rm .} Комплексной аналитической функции $f$ 

%и оператору $\hat A_f$ 

можно сопоставить функцию на расширенном фазовом пространстве \eqref{Psi(qp)}{\rm :}

$$\Psi(q,p)=e^{-\frac{q^2+p^2}4}\,f\Bigl(\frac{{q}-i{p}}{\sqrt2}\Bigr).$$ 

%и оператор $\hat A_\Psi=e^{-\frac{q^2+p^2}4}\hat A_f e^{\frac{q^2+p^2}4}$. 

% Функция $\Psi(q,p)$ задаёт разложение волновой функции из $\Phi'$ по когерентным состояниям. Поскольку когерентные состояния %определены с точностью до фазового множителя, 

Локальные калибровочные преобразования с группой $U(1)$ для функции \emph{на расширенном фазовом пространстве} 

имеют вид $\Psi(q,p)\to\Psi(q,p)\cdot e^{i\alpha(q,p)}.$

%$\hat A_\Psi\to e^{i\alpha(q,p)}\hat A_\Psi e^{-i\alpha(q,p)}$.

\end{newthm}



{\thanks{Работа выполнена при поддержке РФФИ $($проект \rm \No~$11{-}01{-}00828{-}a).$}}



%% Благодарности, гранты и прочее



\begin{thebibliography}{9}



\Bibitem{ivanovmg:bib:Weyl}

\by H.~Weyl

\book The Theory of Groups and Quantum Mechanics

\publ Dover

\publaddr New York

\totalpages 444

\yr 1931

\rtransl русск. пер.:~

\by Г.~Вейль

\book Теория групп и квантовая механика

\publaddr М.

\publ Наука. 

\yr 1986

\totalpages 496



\RBibitem{ivanovmg:bib:Faddeev}

\by Л.~Д.~Фаддеев, О.~А.~Якубовский

\book Лекции по квантовой механике для студентов"=математиков

\publaddr М., Ижевск

\publ РХД

\yr 2001

\totalpages 256

\transl англ. пер.:~

\book Lectures on Quantum Mechanics for Mathematics Students

\by L.~D.~Faddeev, O.~A.~Yakubovskii

\yr 2009

\serial Student Mathematical Library

\publ American Mathematical Society

\vol 47

\publaddr Providence, RI

\totalpages xii+234



\Bibitem{ivanovmg:bib:bellissard}

\paper Noncommutative Geometry of Aperiodic Solids 

\by   J.~Bellissard

\inbook Geometric and Topolo\-gical Methods for Quantum Field Theory

\procinfo Villa de Leyva, 2001

\publ World Sci. Publishing

\publaddr River Edge, NJ

\yr 2003

\pages 86--156

\crossref{http://dx.doi.org/10.1142/9789812705068_0002}

\mathscinet{http://www.ams.org/mathscinet-getitem?mr=2009996}

\zmath{http://www.zentralblatt-math.org/zmath/search/?an=Zbl 1055.81034}



\Bibitem{ivanovmg:bib:deGosson}

\paper A gauge theory of quantum mechanics

\by   J.~M.~Isidro, M.~A.~de~Gosson

\jour Mod. Phys. Lett. A.

\yr 2007

\vol 22

\issue 3

\pages 191--200



\RBibitem{ivanovmg:bib:ivanovmg}

\by М.~Г.~Иванов

\book Как понимать квантовую механику

\publaddr М., Ижевск

\publ РХД

\yr 2012

\totalpages 516

\misctransl

\by M.~G.~Ivanov

\book How to understand quantum mechanics

\publaddr Moscow, Izhevsk

\publ RCD

\yr 2012

\totalpages 516





\end{thebibliography}



\noindent

\hskip6.5cm \thedate \par\noindent

\hskip6.5cm\therevdate



%

%\newpage

%

%\chead[\fancyplain{}{\scriptsize \sl I\,v\,a\,n\,o\,v~M.\,G.}]{\fancyplain{}{\scriptsize \sl  Phase space curvature}}



\makeengtitle







\newpage



%

% \tableofengcontents

%%

%\end{document}